% Section 4: Analysis
% Mechanistic analysis and deeper investigation

\section{Analysis}
\label{sec:analysis}

% ============================================================
% 4.1 Correlation ≠ Transferability
% ============================================================
\subsection{Correlation $\neq$ Transferability}

A striking finding emerges from comparing geometric alignment
with functional transferability.

\textbf{CCA results.}
Canonical Correlation Analysis (CCA) reveals that
first canonical correlation $\rho_1 \approx 1.0$ for both K and V
across all model pairs (Table~\ref{tab:cca}).
This indicates near-perfect linear alignability:
teacher and student KV representations lie in essentially
the same linear subspace.

\begin{table}[t]
\centering
\caption{CCA $\rho_1$ vs functional transferability.
High correlation does not predict transfer success.}
\label{tab:cca}
\begin{tabular}{lrrr}
\toprule
Pair & $\rho_K$ & $\rho_V$ & K-only $\dll$ \\
\midrule
8B$\to$0.6B & 1.000 & 0.9998 & $+4.98$ \\
4B$\to$1.7B & 1.000 & 0.9999 & $+7.31$ \\
4B$\to$0.6B & 1.000 & 0.9999 & $+3.26$ \\
\bottomrule
\end{tabular}
\end{table}

\textbf{The dissociation.}
Despite near-perfect geometric alignment for both K and V,
only K transfers functionally.
V-only $\dll$ is weak or negative (Table~\ref{tab:main}).
This dissociation is the paper's central finding.

\textbf{Interpretation.}
Linear-geometric compatibility (CCA $\rho \approx 1$)
is necessary but not sufficient for functional transfer.
The bottleneck lies in downstream decode,
where content is consumed through weight-parameterized pathways.
K geometry is model-agnostic; V content is weight-bound.

% ============================================================
% 4.2 Layer Localization
% ============================================================
\subsection{Layer Localization of V Advantage}

Per-layer V-only replacement reveals concentrated hotspots
at layers 8 and 12 (Table~\ref{tab:layers}).

\begin{table}[t]
\centering
\caption{Per-layer V-only $\dll$ on 8B$\to$0.6B.
V advantage concentrates at layers 8/12.}
\label{tab:layers}
\begin{tabular}{lr}
\toprule
Layer & V-only $\dll$ \\
\midrule
Layer 4 & $-0.23$ \\
Layer 8 & $+1.07$ \\
Layer 12 & $+3.72$ \\
Layer 16 & $-0.45$ \\
Layer 20 & $-0.89$ \\
Layer 24 & $-1.23$ \\
Layer 28 & $-1.56$ \\
\bottomrule
\end{tabular}
\end{table}

\textbf{Concentration.}
Layer 12 alone accounts for $+3.72$ $\dll$.
Layer 8 contributes $+1.07$.
All other layers show near-zero or negative effects.

\textbf{Asymmetry with K.}
K advantage distributes globally:
no single layer dominates, but global K transfer yields large gains ($+4.98$).
This suggests K geometry is a property of the entire network,
while V content is localized to specific processing stages.

% ============================================================
% 4.3 Cost Analysis
% ============================================================
\subsection{Cost Crossover}

We analyze the computational cost of state transfer vs.\ weight transfer.

\textbf{Mapper cost.}
Affine mapper parameters: 58.78M (224.2 MB at float32).
This cost is fixed regardless of sequence length.

\textbf{Cache cost.}
Per-token KV cache: 112 KB (28-layer student).
This cost scales linearly with sequence length.

\textbf{Crossover point.}
State/weight transfer cost crossover at $\sim$2050 tokens (Figure~\ref{fig:cost}).
Below 2050 tokens, transferring KV cache is cheaper.
Above 2050 tokens, transferring mapper parameters is cheaper.

\begin{figure}[t]
\centering
% Placeholder for cost curve
\caption{State/weight transfer cost ratio vs sequence length.
Crossover at $\sim$2050 tokens.}
\label{fig:cost}
\end{figure}

\textbf{Implication.}
For short-context applications (chat, few-shot QA),
KV state transfer is cost-effective.
For long-context applications (document summarization),
mapper transfer is preferred.

% ============================================================
% 4.4 Boundary Conditions
% ============================================================
\subsection{Boundary Conditions}

Two pairs require special attention:

\textbf{1.7B$\to$0.6B (failure).}
K-only $\dll = -0.91$ (negative).
Under ``K is shared geometry'' theory, this should be neutral.
The negative value suggests weak teacher K can actively hurt.
This establishes a lower bound on teacher capacity:
teacher must exceed student for useful K transfer.

\textbf{8B$\to$4B (saturation).}
Self EM = 100\%, so $\dll = +7.67$ is pure calibration gain,
not task improvement.
The student already solves the task; K transfer only
improves confidence calibration.

% ============================================================
% 4.5 Synthesis
% ============================================================
\subsection{Synthesis}

The evidence converges on a clear picture:

\textbf{K is addressing geometry.}
K states encode attention routing patterns that are
model-agnostic and composable across architectures.
This explains why K transfers: the routing logic is generic.

\textbf{V is weight-bound content.}
V states encode content representations that are
consumed through weight-parameterized pathways.
This explains why V doesn't transfer: the content requires
compatible weights to be useful.

\textbf{Implication.}
For KV cache transfer:
transfer K, not V.
For distillation:
focus on V/content pathways, not addressing geometry.
For model stitching:
addressing geometry is composable; content is not.
